\section{The retained extremal inputs}
The classification proved in this section is Cameron--Cara, Theorem 2.1, with its strong intersection and relative-subgroup consequences. Its proof is now common to two applications. The ordinary and parity-modified families remain their source results; the present organization makes no new historical classification claim.
\textbf{W1 (Whiston bound and extremal generation).} An independent family in \(S_m\) has at most \(m-1\) members, and an independent family with \(m-1\) members generates \(S_m\). We use precisely this implication, not a converse asserting that every independent generating family has maximum size.

\textbf{W2 (Whiston omission consequence).} If \(n\geq7\) and \((s_i)\) is an independent generating family of \(S_n\) of size \(n-1\), each omission subgroup \(\langle s_j:j\ne i\rangle\) is intransitive, is transitive imprimitive with blocks of size two, or is \(A_n\). This is the separate consequence of Whiston's proof stated on page 4 of the 2002 manuscript. Its dependence on the classification of finite simple groups is retained. W2 is not deduced here from W1.
Cayley's labelled-tree count and the innerness of every automorphism of \(S_n\) for \(n\geq7\) are retained classical inputs. The CFSG-dependent omission consequence W2 is separate from W1.

\section{Tree markings and the parity construction}
\textbf{Theorem UB.Sign (sign-kernel extension).} Let the blocks be singletons and distinguish an edge \(d\), with \(s=\tau_d\). For \(e\ne d\), choose \(\epsilon_e\in\{1,-1\}\) and put \begin{equation}
r_d=s,\qquad r_e=(s\tau_e)^{\epsilon_e}.
\label{eq:maximal-symmetric-markings:2-4}
\end{equation} For \(U\subseteq E\), set \(F(U)=\{d\}\cup(U\setminus\{d\})\). Then \begin{equation}
\langle r_e:e\in U\rangle=
\begin{cases}
K_{F(U)},&d\in U,\\
K_{F(U)}\cap A_\Omega,&d\notin U.
\end{cases}
\label{eq:maximal-symmetric-markings:2-5}
\end{equation} Both the ordinary marking \((\tau_e)\) and \eqref{eq:maximal-symmetric-markings:2-4} generate \(S_\Omega\), are independent, and satisfy the full subset intersection property with \(B=1\).

\emph{Proof.} Inverting an individual generator leaves its cyclic subgroup unchanged. If \(s\) is included, \(s\) and \(s\tau_e\) recover \(\tau_e\), proving the first case. The sign kernel of \(K_F\), for \(d\in F\), is generated by the \(s\tau_e\), \(e\in F\setminus\{d\}\): every even word in edge transpositions is a product of pairs, and \[
\tau_e\tau_f=(s\tau_e)^{-1}(s\tau_f),\qquad s\tau_d=1.
\] This proves the second case, including the empty set, since \(K_{\{d\}}\cap A_\Omega=1\). Intersect \eqref{eq:maximal-symmetric-markings:2-5} using Theorem UB.F and the sign condition. Since \(F(U)\cap F(V)=F(U\cap V)\), the result is exactly the subgroup for \(U\cap V\). The same conclusion for ordinary markings follows from the unit-weight specialization of Theorem UB.F. Omitting an ordinary edge or a twisted label \(e\ne d\) leaves its cut invariant; the omitted generator crosses that cut. Omitting \(r_d\) leaves only even permutations, whereas \(r_d\) is odd. These observations prove independence and the generation assertions. ∎
\section{The extremal converse}
\textbf{Corollary UB.Blocks (the two-point-block obstruction).}
If $m\geq4$ and $H\leq C_2\wr S_m$ is transitive on $2m$ points,
then $H$ has no independent generating family of size $2m-2$.
\begin{proof}
Let $N$ be the block kernel and $P$ the transitive block image.
If $P<S_m$, Whiston's equality clause gives $\mu(P)\leq m-2$,
and Lemmas~\ref{lem:quotient} and~\ref{lem:binary} give
$\mu(H)\leq2m-4$. If $P=S_m$, their stronger module bound gives
$\mu(H)\leq m+1$. Both quantities are strictly smaller than $2m-2$
for $m\geq4$, including the endpoint $m=4$.
\end{proof}

\textbf{Theorem UB.Extremal (extremal recognition).} Assume \(n\geq7\). Every independent family of \(n-1\) elements of \(S_n\) is an ordinary unit-weight tree marking from Theorem UB.F or a sign-kernel marking \eqref{eq:maximal-symmetric-markings:2-4}. If the family is independent relative to \(B\leq S_n\), then \(B=1\).

\emph{Proof.} By W1 the family generates \(S_n\). For its omission groups \(G_i\), W2 and Corollary UB.Blocks exclude the two-point-block case: in that case \(n=2m\geq8\), and the remaining \(n-2\) generators are an independent generating family for \(G_i\). At most one \(G_i\) is \(A_n\), because its omitted generator must be the unique odd member. Label a possible exception by \(1\).

For each \(i>1\), choose an edge \(e_i=\{x,s_i(x)\}\) crossing two \(G_i\)-orbits. Such an edge exists, since otherwise every generator would preserve those orbits. Choose \(e_1\) from any moved point of \(s_1\). If all the omission groups are intransitive, choose \(e_1\) crossing its own omission orbits as well. For \(j\ne i\), the edge \(e_j\) lies inside a \(G_i\)-orbit because \(s_j\in G_i\). In particular the chosen edges are distinct: a repeated edge with a label \(i>1\) would both cross and stay inside the \(G_i\)-orbits. Thus a cycle cannot contain \(e_i\) for \(i>1\): the remaining edges would give a path within a \(G_i\)-orbit between its endpoints. A cycle cannot consist of \(e_1\) alone. The \(n-1\) edges therefore form a tree on the \(n\) points.

Deleting \(e_i\), for \(i>1\), leaves two connected sets each contained in a \(G_i\)-orbit; intransitivity makes them exactly its two orbits. Hence \(s_j\) preserves this cut whenever \(j\ne i\). Intersect these cut stabilizers using Theorem UB.F. It follows that \(s_1\in K_{\{e_1\}}\), so \(s_1=\tau_{e_1}\), and that \(s_i\in K_{\{e_1,e_i\}}\). If every omission group is intransitive, include the first cut too to obtain \(s_i=\tau_{e_i}\) for every \(i\). Otherwise \(G_1=A_n\), and all other generators are nonidentity even elements of the two-edge group. For adjacent edges that group is \(S_3\), whose nonidentity even elements are the two inverse three-cycles; for disjoint edges it is \(C_2\times C_2\), with a unique nonidentity even element. These are precisely \((s_1\tau_{e_i})^{\pm1}\), giving \eqref{eq:maximal-symmetric-markings:2-4}.

Now suppose the original family is independent relative to \(B\). It is independent in the ordinary sense, so the preceding classification applies. By Theorems UB.F and UB.Sign its omission groups are the tree-cut stabilizers, except that in the sign-kernel case the distinguished omission group is \(A_n\). At each maximal omission group independence forces \(B\leq G_i\); otherwise \(\langle B,G_i\rangle=S_n\). A pendant edge different from the distinguished edge exists and yields a point stabilizer, so \(B\) fixes a point. At a nonmaximal cut, Lemma UB.Overgroups forces \(B\) into the unordered two-block stabilizer. Its elements outside the ordered stabilizer interchange the halves and fix no point. Thus \(B\leq G_i\) here also. Strong independence from Theorems UB.F and UB.Sign gives \(B\leq\bigcap_iG_i=1\). ∎

